\documentclass[10pt,letterpaper]{article}
\usepackage{cornell-notes}

\title{Clause 23.03.03.03: Deduction Guides}
\author{}
\date{}
\setDocTitle{Clause 23.03.03.03: Deduction Guides}
\setDocSubtitle{C++ 2024}
\setDocOwner{C++ 2024 Cornell Notes}
\setDocDate{\today}
\setCornellCollection{C++ 2024 Cornell Notes}
\setCornellUnitType{Clause}
\setCornellUnitNumber{23.03.03.03}
\setCornellUnitTitle{Deduction Guides}
\hypersetup{pdftitle={Clause 23.03.03.03: Deduction Guides},pdfsubject={Cornell notes},pdfauthor={}}

\begin{document}
\maketitle
\begin{CornellOverviewBox}{Overview}
\textbf{Classification:} Standard library - Normative\\[2pt]
\textbf{Learning objective:} Understand the rules, terminology, and semantic consequences associated with Deduction guides.
\end{CornellOverviewBox}\begin{CornellSummaryBox}{Summary}
{\sffamily\bfseries\color{Accent} Summary}\par\vspace{3pt}Section 23.3.3.3 focuses on Deduction guides. Use the listed grammar productions together with the semantic constraints and disambiguation rules referenced by the section. When applying it, preserve the distinction between mandatory requirements, recommendations, permissions, and behavior deliberately left undefined, unspecified, or implementation-defined.
\end{CornellSummaryBox}

\end{document}
